% !TeX program = pdflatex
%% IEEJ English manuscript (ieej-e.cls). Upload this folder to Overleaf as project root.
\documentclass[english,fleqn]{ieej-e}
\usepackage{silence}
\WarningFilter{caption}{Unknown document class (or package)}
\WarningFilter{fixltx2e}{fixltx2e is not required}
\WarningFilter{latex}{Text page 5 contains only floats.}
\usepackage[defaultsups]{newtxtext}
\usepackage[varg]{newtxmath}
\usepackage[superscript,nomove]{cite}
\usepackage{graphicx}
\usepackage{subcaption}
\usepackage{booktabs}
\usepackage{multirow}
\usepackage{array}
\setcounter{topnumber}{6}
\setcounter{bottomnumber}{6}
\setcounter{totalnumber}{12}
\renewcommand{\topfraction}{0.92}
\renewcommand{\bottomfraction}{0.92}
\renewcommand{\textfraction}{0.07}
\renewcommand{\floatpagefraction}{0.72}
\renewcommand{\dbltopfraction}{0.95}
\renewcommand{\dblfloatpagefraction}{0.72}
\setlength{\dbltextfloatsep}{8pt plus 4pt minus 3pt}
\setlength{\dblfloatsep}{8pt plus 4pt minus 2pt}
\usepackage{url}
\usepackage{xurl}
\usepackage[pdfencoding=auto]{hyperref}
\hypersetup{%
 setpagesize=false,
 colorlinks=true,
 urlcolor=blue,
 citecolor=black,
 linkcolor=black,
}

%% --- Journal metadata (fill before submission) ---
\FIELD{}
\YEAR{}
\NO{}
%% \received{2026}{3}{15}
%% \revised{}{}{}

\title{Bias Evaluation of MediaPipe Pose Estimation\\
       Using a Unity Simulation Environment}

\authorlist{%
 \authorentry[Fumimaro Taira, E-mail: fmimaro.work0217@gmail.com]{Fumimaro Taira}{n}{KHS}
 \authorentry{Tsukasa Kato}{m}{UEC}
 \authorentry{Jin Afuso}{n}{LOL}
}
\affiliate[KHS]{Okinawa Prefectural Kaiko High School\\ Japan}
\affiliate[UEC]{University of the Ryukyus\\ 1, Senbaru, Nishihara-cho, Nakagami, Okinawa 902--0213, Japan}
\affiliate[LOL]{lollol Co. Ltd.\\ Office Izumisaki 2-A, 2--3--3, Izumisaki, Naha, Okinawa 900--0021, Japan}

\setlength{\emergencystretch}{3em}
\hbadness=10000
\vbadness=10000
\hfuzz=120pt

\makeatletter
\def\ft@mailaddress#1{%
 \protected@xdef\@mailaddress{%
  \@mailaddress
  \protect\footnotetext[0]{\protect\raggedright #1\par}}%
}
\makeatother

\begin{document}
\begin{abstract}
This study quantifies systematic biases in MediaPipe pose estimation using
ground truth from a Unity simulation with 505 multi-view cameras on four
height layers.
The GT is synthetic; the reported statistics therefore describe estimator
behaviour under these reproducible conditions, and generalisation to real
subjects and in-the-wild video requires independent validation.
The results support \emph{ancillary, research-oriented} use of pose outputs
once biases and estimator limits are acknowledged; they are not a drop-in
replacement for safety-critical pipelines (e.g.\ detection, ranging, time-to-contact).
Whereas prior evaluations emphasise single-number accuracy under limited
viewpoint conditions, we identify structured error patterns arising from model
architecture, training-data priors, and coordinate conventions.
We collected 259{,}356 joint observations (107 frames $\times$ 505 cameras
$\times$ 12 joints) and analysed joint-angle MAE and direction-angle errors
($\Delta\theta$, $\Delta\psi$).
Four systematic biases are identified and linked to primary evidence:
coordinate-system bias mitigated by $y \leftarrow -y$
(Table~\ref{tab:yflip_effect}); a pelvic rotation artefact from
anti-correlated hip $\Delta\psi$ ($r=-0.84$); upper-limb coupling in
$\Delta\psi$ ($r=0.72$--$0.77$); and viewpoint-dependent joint-angle MAE,
with overhead ($Y{=}2.0$\,m) cameras showing up to approximately 15\,\% higher
MAE than at $Y{=}1.0$\,m (Table~\ref{tab:mae_by_layer}), consistent with a
frontal prior in training data (Figure~\ref{fig:mae_layer_compare}).
\end{abstract}
\begin{keyword}
MediaPipe, BlazePose, pose estimation, bias evaluation, Unity simulation,
direction angle error
\end{keyword}
\maketitle
\section{Introduction}

Pose estimation has become a core enabling technology across sports
analysis, clinical rehabilitation, ergonomics, and human--computer
interaction~\cite{mediapipe,blazepose}.
\textbf{Low-latency} estimation is especially important for applications that
must react to streaming video, including sports broadcast analytics, online
training feedback, safety monitoring in factories and warehouses, motion
observation in rehabilitation, and XR or in-game avatar control.
MediaPipe, developed by Google, is particularly attractive for practitioners
because it enables real-time processing on commodity hardware without GPU
acceleration~\cite{mediapipe}.
In surveillance-like or automotive-like settings where viewpoint, lighting,
and occlusion are not controlled, full-body pose is \textbf{unlikely to
replace} primary pipelines based on object detection, ranging, or time-to-contact
estimation; it is more realistically positioned as \textbf{auxiliary
information} or as an \textbf{R\&D evaluation axis} once biases and limits
are understood.
This work also aims to quantify that premise.

Despite its widespread deployment, a systematic and multi-viewpoint
quantitative evaluation of MediaPipe's 3D accuracy remains lacking.
Bazarevsky~et~al.~\cite{blazepose} report BlazePose benchmark performance,
but their evaluation is restricted to a curated benchmark with limited
viewpoint diversity. In real-world deployments, camera placement is
constrained by the environment and cannot be controlled; understanding how
accuracy varies with viewpoint is therefore essential.

Existing independent evaluations of BlazePose and related estimators
(reviewed in detail in Section~2) have focused on accuracy under standard
benchmark conditions and near-frontal views; systematic multi-viewpoint
assessment across diverse camera heights and horizontal positions is
absent. Simulation-based ground truth collection (pioneered by
SURREAL~\cite{surreal} and extended by Unity-based
generators~\cite{peoplesanspeople}) provides a scalable alternative to
physical motion capture, permitting fully controlled camera placement across
arbitrary viewpoints without additional instrumentation.

A fundamental challenge for any monocular 3-D pose estimator is depth
ambiguity: a single 2-D image projection does not uniquely determine 3-D
joint positions, and estimation accuracy depends strongly on the training-data
viewpoint distribution~\cite{mehta2017}. Photo-centric public datasets such as
COCO~\cite{coco} contain diverse poses per instance, yet standing- or
walking-like postures together with near-frontal camera viewpoints are
\emph{relatively} common and are often discussed as biasing learned models
toward frontal viewpoints, which \emph{may} impair generalisation to lateral and
overhead cameras.
We examine this hypothesis across 505 camera positions and additionally
identify a systematic Y-axis coordinate mismatch between Unity and MediaPipe
that, if uncorrected, yields spuriously large mean $|\Delta\theta|$ for
proximal joints (e.g., $\approx121^{\circ}$ for elbows before
$y \leftarrow -y$).

The contributions of this paper are as follows:
\begin{enumerate}\itemsep0pt
  \item \textbf{Bias taxonomy for MediaPipe}: We define four bias classes
        with explicit operational meaning and link each to primary evidence:
        coordinate $y$-sign effects (Table~\ref{tab:yflip_effect};
        Sec.~3.4); pelvic depth anti-correlation
        (Figure~\ref{fig:hip_psi_heatmap}, Table~\ref{tab:correlation_psi});
        upper-limb error coupling
        (Tables~\ref{tab:correlation_theta}--\ref{tab:correlation_psi},
        Figure~\ref{fig:correlation_heatmaps}); and viewpoint-dependent MAE
        (Table~\ref{tab:mae_by_layer}, Figure~\ref{fig:mae_layer_compare}).
  \item \textbf{Coordinate system bias correction}: We identify a
        systematic Y-axis orientation mismatch between Unity and MediaPipe,
        quantify its per-joint impact (up to $161^{\circ}$ reduction in
        $|\Delta\theta|$), and provide the correction formula
        $y \leftarrow -y$.
  \item \textbf{Direction angle metrics}: We introduce $\Delta\theta$ and
        $\Delta\psi$ as bias-sensitive metrics complementary to joint angle
        MAE; these metrics reveal structured error patterns (in particular the
        pelvic artefact) that aggregate MAE conceals.
  \item \textbf{Large-scale multi-viewpoint dataset}: 505 camera positions
        $\times$ 107 frames yield 259,356 observations covering extreme
        lateral and overhead viewpoints absent from standard benchmarks,
        enabling, to our knowledge, the first systematic viewpoint-bias analysis
        for MediaPipe in the open literature.
\end{enumerate}

This paper is organised as follows:
Section~2 reviews related work;
Section~3 describes the experimental method;
Section~4 presents results;
Section~5 discusses implications;
Section~6 concludes.

% ============================================================
% 2. RELATED WORK  (target: ~1 page) --- NEW SECTION
% ============================================================
\section{Related Work}

\subsection{MediaPipe and BlazePose Accuracy}

BlazePose, the neural-network backbone of MediaPipe Pose, was introduced by
Bazarevsky~et~al.~\cite{blazepose}, who report PCK@0.2 of 72.0\,\% on the
COCO validation set and state-of-the-art performance on a proprietary
fitness-video benchmark. Their evaluation is, however, restricted to
near-frontal perspectives typical of fitness recordings; performance under
diverse camera viewpoints is not assessed.

The original BlazePose evaluation~\cite{blazepose} does not provide
systematic accuracy data across diverse camera positions.
The present study addresses this gap with a
505-camera 3-D grid spanning four height layers (Y\,=\,0.5--2.0\,m).

\subsection{Simulation-Based Ground Truth Collection}

Obtaining accurate 3-D ground truth at scale is impractical in real-world
settings: manual 3-D annotation is labour-intensive and optical
motion-capture systems are costly and laboratory-confined. Synthetic
rendering has therefore become a standard approach. Varol~et~al.~\cite{surreal}
introduced SURREAL, rendering over 6~million photo-realistic frames of
SMPL-body avatars with perfect 2-D/3-D pose, depth, and segmentation labels;
models pre-trained on SURREAL transfer effectively to real-world benchmarks.
Patel~et~al.~\cite{agora} extended this
paradigm with AGORA, compositing 4,240 high-quality textured body scans into
natural outdoor environments and demonstrating that training on realistic
synthetic data improves regression performance on the 3DPW benchmark.

Game-engine-based capture systems add the further advantage of full camera
controllability. Ebadi~et~al.~\cite{peoplesanspeople} released
PeopleSansPeople, a Unity-based synthetic data generator that produces
COCO-format keypoint annotations; pre-training on PeopleSansPeople data and
fine-tuning on real COCO images achieves keypoint AP of 63.5, surpassing a
real-data-only baseline of 62.0~AP in the high-data regime. Our study
similarly exploits a Unity simulation environment to record exact per-frame
3-D joint coordinates from 505 systematically placed cameras (a viewpoint
sweep that physical motion-capture installations cannot economically
reproduce).

\subsection{Monocular 3D Pose Estimation and Depth Ambiguity}

Recovering 3-D pose from a single RGB camera is inherently ill-posed: many
3-D configurations project to the same 2-D image under perspective
projection. Mehta~et~al.~\cite{mehta2017} improved generalisation to
in-the-wild scenes through multi-dataset transfer learning and introduced the
MPI-INF-3DHP benchmark, which includes greater viewpoint diversity than
standard pose benchmarks; nevertheless, lateral and overhead perspectives
remain under-represented in training corpora.

Public datasets such as COCO~\cite{coco} are widely discussed as skewing
toward near-frontal viewpoints and standing- or walking-like poses, which
\emph{may} embed a frontal prior in learned models and \emph{may} reduce
accuracy for lateral or overhead cameras; our multi-height evaluation
quantifies this tendency (Table~\ref{tab:mae_by_layer}).
Beyond viewpoint bias, MediaPipe's skeleton model represents the pelvis as a
simple line segment connecting landmarks~23 (left hip) and~24 (right hip),
without a rigid multi-segment structure. A small horizontal pixel offset in
the 2-D image can therefore be misinterpreted as a large pelvic rotation
during 3-D lifting, producing the symmetric anti-correlated hip depth error
($r=-0.840$) reported in Section~4.

\subsection{Skeletal Constraint Models}

Several architectures impose anatomical or temporal constraints to mitigate
depth ambiguity. Mehta~et~al.~\cite{vnect} demonstrated real-time monocular
3-D pose estimation with a kinematic skeleton fitting stage that enforces
bone-length consistency and temporal stability across frames.
PoseFormer~\cite{poseformer} is a purely transformer-based video pose
estimator that captures long-range spatio-temporal joint relations without
convolutional layers, achieving 44.3\,mm MPJPE on Human3.6M.
MotionBERT~\cite{motionbert} pre-trains a Dual-stream Spatio-temporal
Transformer to recover 3-D motion from noisy 2-D observations, embedding
geometric and physical priors that generalise to mesh recovery and action
recognition, achieving 39.2\,mm MPJPE on Human3.6M.

Despite these advances, pelvic rigidity is rarely enforced as an explicit
hard constraint: the pelvis is typically modelled as a single segment or
handled implicitly through temporal consistency. The hip anti-correlation
($r=-0.840$) observed in our data suggests that constraining left- and
right-hip depth to be mutually consistent under a rigid pelvic model could
substantially reduce depth errors for walking motions (a direction we
recommend for future model design).

% ============================================================
% 3. EXPERIMENTAL METHOD  (target: ~1.5 pages)
% ============================================================
\section{Experimental Method}

\subsection{Experimental Environment}

Using Unity (version 6000.0.60f1), we constructed a walking simulation
environment with a Y-Bot humanoid avatar. The avatar performed a straight
walking motion from coordinates $(0,\,0,\,-3)$ to $(0,\,0,\,3)$ along the
positive Z-axis.

\begin{figure}[!t]
\centering
\includegraphics[width=\linewidth]{figs/fig01_camera_layout.png}
\caption{Camera layout (505 positions across four height layers Y = 0.5,
  1.0, 1.5, 2.0). The centre arrow indicates the avatar's walking direction
  (+Z axis).}
\label{fig:camera_layout}
\end{figure}

Cameras were placed on a 3-D grid covering four height layers
(Y = 0.5, 1.0, 1.5, 2.0\,m) and horizontal positions
(X, Z $\in [-6, 6]$\,m). Within each layer a $13\times13$ grid was used
with the central $5\times5$ block removed, giving a theoretical maximum of
$4(13^2-5^2)=576$ positions. Data were collected at 505 of these positions.
Each position yielded 107 captured frames at a resolution of
$1280\times720$\,pixels, for a total of 259,356 joint observations
(107 frames $\times$ 505 cameras $\times$ 12 joints, after visibility
filtering). The same GT and image data are reusable for evaluating
alternative estimators such as YOLO-Pose.

The humanoid avatar was a Y-Bot mesh (Unity standard humanoid rig) driven
by a built-in looped walking animation clip. The motion spans approximately
6\,m along the positive Z-axis and covers a full walking cycle, providing
natural joint-angle variation across frames.

% NOTE: fig02a (Unity実験環境.png) and fig02b (frame_0018.jpg) are sourced
%   from the original submitted paper (paper/submitted/Zeval_IEEJ_submit.pdf).
%   Current placeholders in paper/source/figs/ should be replaced with the
%   extracted originals before final Overleaf upload.
\begin{figure}[!t]
  \centering
  \includegraphics[width=\linewidth]{figs/fig02a_unity_env.png}
  \vspace{0.3em}
  \includegraphics[width=\linewidth]{figs/fig02b_frame_sample.jpg}
  \caption{Unity experiment environment (top) and a representative captured
    frame (bottom; camera at $Y{=}0.5$\,m, $X{=}{-}1.0$\,m, $Z{=}{-}3.0$\,m).}
  \label{fig:unity_env}
\end{figure}

\subsection{Capture System}

A Trigger Zone Capture System was implemented in Unity.
\texttt{CaptureSystemManager} handled character generation, animation
playback, per-frame image capture, and joint coordinate recording.
\texttt{AutoCaptureManager} read the 505 camera positions from a CSV file
and triggered automatic capture at each position. This ensured consistent,
reproducible multi-viewpoint data collection with a single character motion
sequence.

\subsection{Data Collection}

The Unity simulation ran at 30\,fps; 107 frames were uniformly sampled from
the walking cycle for analysis.
Joint coordinates (GT) were exported per frame in CSV format using Unity's
\texttt{HumanBodyBones} API~\cite{unity}.
MediaPipe's Pose Landmarker (Python package \texttt{mediapipe} 0.10,
Solutions API; \texttt{model\_complexity=1},
\texttt{static\_image\_mode=True},
\texttt{min\_detection\_confidence=0.5}) was applied to each saved image to
extract 3-D coordinates for 33 landmarks. Landmarks with a visibility score
below 0.5 were excluded from analysis to limit the effect of heavy
occlusion. The 12 joints retained for analysis are: left/right shoulder,
elbow, wrist, hip, knee, and ankle.

\subsection{Angle Computation and Error Evaluation}

We report \textbf{two complementary error metrics}.

\subsubsection{Joint Angle Error (Joint Angle MAE)}

The flexion angle at each joint was computed from three anatomical points:
upper arm--elbow--forearm for the elbow; thigh--knee--lower leg for the knee;
and equivalent triplets for shoulder and hip. The angle $\alpha$ at the
vertex between the two vectors was computed as
\begin{equation}
  \alpha = \arccos\!\left(\frac{\mathbf{u}\cdot\mathbf{v}}
                               {\|\mathbf{u}\|\,\|\mathbf{v}\|}\right),
  \quad \alpha \in [0^{\circ},\,180^{\circ}].
  \label{eq:joint_angle}
\end{equation}
The mean absolute difference between GT and MediaPipe angles over all valid
frames and cameras gives the MAE for each joint.

\subsubsection{Direction Angle Error ($\Delta\theta$, $\Delta\psi$)}

Following the BlazePose output specification~\cite{blazepose}, relative
joint positions are expressed with the hip midpoint as origin. In this
hip-centred coordinate system we define
\begin{equation}
  \theta = \mathrm{arctan2}(y,\,x), \quad
  \psi   = \mathrm{arctan2}(z,\,x),
  \label{eq:direction_angles}
\end{equation}
where $(x,y,z)$ is the joint position relative to the hip centre.
The angular errors $\Delta\theta$ and $\Delta\psi$ are the signed
differences between GT and MediaPipe values, normalised to
$(-180^{\circ}, +180^{\circ}]$.

\subsubsection{Coordinate System Unification}

Unity uses a left-handed coordinate system with Y-axis pointing
upward, whereas MediaPipe's \texttt{pose\_landmarks} uses an image
coordinate system with Y-axis pointing downward. This discrepancy
introduces a systematic sign flip in $\theta = \mathrm{arctan2}(y,x)$.
We therefore negate the MediaPipe $y$ coordinate (in practice,
$y \leftarrow -y$) to align with Unity's upward $Y$.
The procedure is implemented in
\path{05_direction_detection/scripts/coordinate_transform.py}.
All direction angle results and inter-joint correlation analyses reported
below use coordinates after this unification.

\subsubsection{Reproducibility and pipeline paths}

Table~\ref{tab:repro_pipeline} lists the main repository paths used to
regenerate figures and statistics. MediaPipe inference used
\path{model_complexity=1}, \path{static_image_mode=True},
\path{min_detection_confidence=0.5}; landmarks with visibility below $0.5$
were dropped before joint-wise analysis. Ground-truth joint positions were
exported from Unity (\path{HumanBodyBones}). Joint-angle MAE aggregates are
obtained by concatenating \path{coordinate_angle_mae.csv} from
\path{03_joint_angle_mae/Y=0.5/} through \path{Y=2.0/}; direction-angle
summaries and correlations use
\path{05_direction_detection/output/processed_data/} and
\path{05_direction_detection/output/correlation_analysis/}
(\path{joint_summary.csv}, \path{detailed_results.csv}, etc.); per-frame
best/worst-camera statistics come from \path{frame_camera_summary.csv}.

\begin{table}[!t]
\centering
\caption{Primary scripts and CSV outputs for reproducing reported numbers
  (paths relative to the analysis repository root).}
\label{tab:repro_pipeline}
\footnotesize
\begin{tabular}{>{\raggedright\arraybackslash}p{0.28\linewidth}>{\raggedright\arraybackslash}p{0.62\linewidth}}
\toprule
Stage & Path \\
\midrule
Y-axis unification & \path{05_direction_detection/scripts/coordinate_transform.py} \\
Joint-angle MAE & \path{03_joint_angle_mae/Y=0.5/}--\path{Y=2.0/} (each \path{coordinate_angle_mae.csv}, four layers) \\
Direction-angle stats & \path{05_direction_detection/output/processed_data/joint_summary.csv} \\
Hip detailed stats & \path{05_direction_detection/output/processed_data/detailed_results.csv} \\
Correlation tables & \path{05_direction_detection/output/correlation_analysis/*.csv} \\
Camera best/worst $|\Delta\theta|$ & \path{05_direction_detection/output/processed_data/frame_camera_summary.csv} \\
\bottomrule
\end{tabular}
\end{table}

\noindent\textbf{Public repositories.}
The Unity capture project is at
\url{https://github.com/MASHMARO2718/bias_evaluation_mediapipe_unity_enenvironment};
analysis scripts for the numbers and figures in this paper are at
\url{https://github.com/MASHMARO2718/bias_evalutation_mediapipe_unity_ieej2026}
(Table~\ref{tab:repro_pipeline} uses paths under the latter root).
Processed MediaPipe \path{CapturedFrames_*.csv} archives (expand under
\path{02_mediapipe_processed/}) are cited as~\cite{zenodo_mp_csv}.

\subsubsection{Statistical analysis}

Pearson correlation coefficients $r$ in Tables~\ref{tab:correlation_theta}
and~\ref{tab:correlation_psi} summarise \emph{linear} co-variation between
paired columns of pooled errors (e.g.\ $\Delta\psi_{\mathrm{L}}$ and
$\Delta\psi_{\mathrm{R}}$ for the same frame and camera). For paired
observations $(x_i,y_i)$, $i=1,\ldots,n$,
\begin{equation}
  r=\frac{\sum_{i=1}^{n}(x_i-\bar{x})(y_i-\bar{y})}
         {\sqrt{\sum_{i=1}^{n}(x_i-\bar{x})^2}\,\sqrt{\sum_{i=1}^{n}(y_i-\bar{y})^2}},
  \qquad -1\le r\le 1.
  \label{eq:pearson_r}
\end{equation}
Here $n$ is the number of pooled rows (each joint pair satisfies
$n>21{,}000$). Two-sided tests of $H_0:\rho=0$ are summarised as
$p<0.001$ in the table captions.

The \textbf{SD} in Table~\ref{tab:direction_angle_error} is the sample
standard deviation of one error type for one joint over $n$ pooled
observations $e_1,\ldots,e_n$:
$\mathrm{SD}=\sqrt{\frac{1}{n-1}\sum_{i=1}^{n}(e_i-\bar{e})^2}$
(means and SDs are computed per joint; coupling between joints is reported
separately via~$r$).

Approximate 95\% confidence intervals for the population correlation
$\rho$ use Fisher's $z=\mathrm{arctanh}(r)$ with
$\mathrm{SE}(z)\approx 1/\sqrt{n-3}$ (large-sample normal approximation for
$z$ with mean $\mathrm{arctanh}(\rho)$ and variance $1/(n-3)$). This
\textbf{SE} gauges uncertainty in inferring $\rho$ from the sample~$r$ and
differs from the SD of angular errors above. Because $n$ is large, intervals
are narrow even for $|r|\approx 0.8$ (Table~\ref{tab:bias_claim_evidence}).

For joint-angle MAE across camera height layers
(Table~\ref{tab:mae_by_layer}), we report descriptive means per layer;
formal omnibus tests (e.g., repeated-measures or clustered non-parametric
comparisons across layers) are not central to the bias-characterisation
goal here but should accompany any claim of ``statistically significant
layer differences'' in follow-on work.

\begin{figure}[!t]
\centering
\includegraphics[width=\linewidth]{figs/fig03_concept_metrics.png}
\caption{\emph{Simplified} schematic of the two evaluation metrics (for
  geometry intuition only). Left: joint angle (MAE); flexion from three
  points ($0$--$180^{\circ}$). Right: direction angles ($\Delta\theta$,
  $\Delta\psi$); offsets relative to the hip centre
  ($-180^{\circ}$ to $+180^{\circ}$).}
\label{fig:concept_metrics}
\end{figure}

% ============================================================
% 4. RESULTS  (target: ~2.5 pages)
% ============================================================
\section{Results}

\subsection{Bias claims, evidence, and effect sizes}

Table~\ref{tab:bias_claim_evidence} links each qualitative bias claim to its
primary figure/table, an illustrative numeric summary, and statistical notes
(correlations: two-sided $p<0.001$ vs.\ $H_0\!:\rho=0$; 95\% CI for $\rho$ via
Fisher $z$ with $n\approx 2.2\times 10^{4}$ observations per pair, yielding
width $\ll 0.02$ for $|r|\approx 0.8$).

\begin{table*}[!t]
\centering
\caption{Bias claims mapped to evidence, illustrative effect sizes, and
  statistical notes ($n\approx 2.2\times 10^{4}$ per correlation pair).}
\label{tab:bias_claim_evidence}
\footnotesize
\begin{tabular}{>{\raggedright\arraybackslash}p{0.19\textwidth} >{\raggedright\arraybackslash}p{0.21\textwidth} >{\raggedright\arraybackslash}p{0.26\textwidth} >{\raggedright\arraybackslash}p{0.28\textwidth}}
\toprule
Bias claim & Primary evidence & Illustrative effect size & Significance / CI \\
\midrule
Coordinate $y$-sign (pre-unification) &
Tab.~\ref{tab:yflip_effect}; Sec.~3.4 &
Shoulder mean $|\Delta\theta|$ reduction
$\approx\!161^{\circ}$ (right shoulder row) &
Deterministic transform; not a stochastic hypothesis test \\
Pelvic / hip depth coupling &
Tab.~\ref{tab:correlation_psi}; Fig.~\ref{fig:hip_psi_heatmap} &
LEFT\_HIP--RIGHT\_HIP: $r=-0.840$ &
$p<0.001$; 95\% CI $\approx[-0.844,-0.836]$ \\
Upper-limb coupling ($\Delta\psi$) &
Tabs.~\ref{tab:correlation_theta}, \ref{tab:correlation_psi};
Fig.~\ref{fig:correlation_heatmaps} &
RIGHT\_ELBOW--RIGHT\_SHOULDER: $r=0.770$ (example) &
$p<0.001$; 95\% CI $\approx[\,0.763,\,0.777\,]$ \\
Viewpoint / overhead MAE &
Tab.~\ref{tab:mae_by_layer}; Fig.~\ref{fig:mae_layer_compare} &
Right hip MAE $+4.7^{\circ}$ vs.\ Y=1.0 row ($+15\%$) &
Descriptive layer means; inferential layer tests optional \\
Best vs.\ worst camera ($|\Delta\theta|$) &
Sec.~\ref{subsec:camera_dep} &
Mean ratio $\approx 9\times$ ($8.8^{\circ}$ vs.\ $80.1^{\circ}$) &
Descriptive per-frame camera ranking \\
\bottomrule
\end{tabular}
\end{table*}

\subsection{Joint Angle MAE}

Table~\ref{tab:joint_angle_error} summarises MAE, median, and maximum
error for eight major joints. Elbow and knee errors are lowest and shoulder
MAE is highest, consistent with greater depth ambiguity for proximal joints;
numerical detail is left to the table.

\begin{table}[!t]
\centering
\caption{Statistics of joint angle error (Joint Angle MAE, degrees).
  Source: \texttt{coordinate\_angle\_mae.csv} (under \texttt{03\_joint\_angle\_mae},
  combined across four camera height layers).}
\label{tab:joint_angle_error}
\begin{tabular}{lccc}
\toprule
Joint          & MAE    & Median & Max    \\
               & (deg)  & (deg)  & (deg)  \\
\midrule
Left shoulder  & 39.8   & 39.2   & 70.1   \\
Right shoulder & 39.0   & 40.0   & 58.2   \\
Left elbow     & 18.6   & 15.9   & 56.9   \\
Right elbow    & 18.2   & 13.6   & 46.2   \\
Left hip       & 30.2   & 29.4   & 66.0   \\
Right hip      & 33.3   & 32.2   & 81.7   \\
Left knee      & 17.0   & 16.6   & 43.4   \\
Right knee     & 19.5   & 19.4   & 62.3   \\
\bottomrule
\end{tabular}
\end{table}

A breakdown of joint angle MAE by camera height layer is provided in
Table~\ref{tab:mae_by_layer} and Figure~\ref{fig:mae_layer_compare};
errors increase noticeably at Y\,=\,2.0\,m (overhead), consistent with
the frontal-view prior embedded in BlazePose training data.

\subsection{Direction Angle Error ($\Delta\theta$, $\Delta\psi$)}

Table~\ref{tab:direction_angle_error} presents the direction angle error
statistics for all 12 joints after coordinate system unification.

\begin{table}[!t]
\centering
\caption{Direction angle error statistics after coordinate unification
  (degrees; absolute-value means). Source: processed joint-summary CSV.
  $^\ddagger$Hip values are computed from detailed results
  ($n=21{,}613$ per hip joint) because hip rows are omitted from the joint summary.}
\label{tab:direction_angle_error}
\small
\begin{tabular}{lcccc}
\toprule
Joint & \multicolumn{2}{c}{$|\Delta\theta|$} & \multicolumn{2}{c}{$|\Delta\psi|$} \\
\cmidrule(lr){2-3}\cmidrule(lr){4-5}
      & Mean & SD & Mean & SD \\
\midrule
Left shoulder  & 10.4 & 11.7 & 94.1 & 86.4 \\
Right shoulder &  9.3 & 10.1 & 92.7 & 92.0 \\
Left elbow     & 58.3 & 24.1 & 87.3 & 90.6 \\
Right elbow    & 59.3 & 22.7 & 88.6 & 94.5 \\
Left wrist     & 84.3 & 103.7 & 88.6 & 91.9 \\
Right wrist    & 91.5 & 111.9 & 91.6 & 99.2 \\
Left hip$^\ddagger$  & 83.5 & 72.5 & 88.9 & 26.4 \\
Right hip$^\ddagger$ & 96.5 & 72.5 & 91.1 & 26.4 \\
Left knee      & 17.4 & 18.8 & 84.7 & 97.9 \\
Right knee     & 18.5 & 19.1 & 87.2 & 97.5 \\
Left ankle     & 11.3 & 14.0 & 90.6 & 105.8 \\
Right ankle    & 11.9 & 14.5 & 100.4 & 103.1 \\
\bottomrule
\end{tabular}
\end{table}

Figure~\ref{fig:elbow_heatmaps} shows the camera-position-wise distribution
of elbow $\Delta\theta$.

\begin{figure}[!t]
  \centering
  \begin{subfigure}[b]{0.48\linewidth}
    \includegraphics[width=\linewidth]%
      {figs/fig04a_elbow_L_theta_y05.jpg}
    \caption{Left elbow $\Delta\theta$ (Y=0.5)}
  \end{subfigure}\hfill
  \begin{subfigure}[b]{0.48\linewidth}
    \includegraphics[width=\linewidth]%
      {figs/fig04b_elbow_R_theta_y05.jpg}
    \caption{Right elbow $\Delta\theta$ (Y=0.5)}
  \end{subfigure}
  \caption{Heatmaps of elbow direction angle error $\Delta\theta$ (Y=0.5).
    Left and right elbows show opposite-sign patterns, consistent with the
    strong negative inter-joint correlation ($r=-0.862$).}
  \label{fig:elbow_heatmaps}
\end{figure}

Figure~\ref{fig:hip_psi_heatmap} shows the hip $\Delta\psi$ heatmap.
Left and right hip exhibit near-identical absolute-value patterns due to
their strong negative correlation ($r=-0.84$).

\begin{figure}[!t]
  \centering
  \begin{subfigure}[b]{0.48\linewidth}
    \includegraphics[width=\linewidth]{figs/fig05a_hip_L_psi_y05.jpg}
    \caption{Left hip $\Delta\psi$ (Y=0.5)}
  \end{subfigure}\hfill
  \begin{subfigure}[b]{0.48\linewidth}
    \includegraphics[width=\linewidth]{figs/fig05b_hip_R_psi_y05.jpg}
    \caption{Right hip $\Delta\psi$ (Y=0.5)}
  \end{subfigure}
  \caption{Heatmaps of hip direction angle error $\Delta\psi$ (Y=0.5).
    Left and right hips show near-identical absolute-value patterns
    because $\Delta\psi_{\text{L}} \approx -\Delta\psi_{\text{R}}$
    ($r = -0.840$): when one hip is forward, the other is backward.}
  \label{fig:hip_psi_heatmap}
\end{figure}

Table~\ref{tab:mae_by_layer} breaks down joint angle MAE by camera height
layer. Errors are generally lowest at Y=1.0 and increase at Y=2.0
(overhead view), suggesting that top-down perspectives reduce estimation
accuracy, consistent with underrepresentation of such viewpoints in training
data. Figure~\ref{fig:mae_layer_bar} summarises this trend across all joints,
and Figure~\ref{fig:mae_layer_compare} shows heatmaps for the left
elbow across three height layers.

\begin{figure*}[!tb]
  \centering
  \begin{subfigure}[b]{0.31\textwidth}
    \includegraphics[width=\linewidth]{figs/fig07a_mae_elbow_L_y05.png}
    \caption{Y=0.5}
  \end{subfigure}\hfill
  \begin{subfigure}[b]{0.31\textwidth}
    \includegraphics[width=\linewidth]{figs/fig07b_mae_elbow_L_y15.png}
    \caption{Y=1.5}
  \end{subfigure}\hfill
  \begin{subfigure}[b]{0.31\textwidth}
    \includegraphics[width=\linewidth]{figs/fig07c_mae_elbow_L_y20.png}
    \caption{Y=2.0}
  \end{subfigure}
  \caption{Joint angle MAE heatmaps for the left elbow at three camera
    height layers. Error increases at Y=2.0 (overhead), consistent with
    training-data viewpoint bias. Source:
    \texttt{03\_joint\_angle\_mae/} (per-layer \texttt{coordinate\_angle\_mae.csv}).}
  \label{fig:mae_layer_compare}
\end{figure*}

\begin{figure}[!t]
  \centering
  \includegraphics[width=\linewidth]{figs/fig08_mae_layer_bar.png}
  \caption{Joint angle MAE by camera height layer (Y = 0.5, 1.0, 1.5, 2.0\,m).
    Error increases consistently at Y\,=\,2.0 for all joints, with the largest
    degradation at the hip (\(+4.7^{\circ}\) for right hip vs.\ Y\,=\,1.0).
    Source: \texttt{03\_joint\_angle\_mae/} (per-layer \texttt{coordinate\_angle\_mae.csv}).}
  \label{fig:mae_layer_bar}
\end{figure}

\begin{table}[!t]
\centering
\caption{Joint angle MAE (deg) by camera height layer.
  Source: \texttt{03\_joint\_angle\_mae/Y=0.5/}--\texttt{Y=2.0/}
  \texttt{coordinate\_angle\_mae.csv} (four layers).}
\label{tab:mae_by_layer}
\small
\begin{tabular}{lcccc}
\toprule
Joint & Y=0.5 & Y=1.0 & Y=1.5 & Y=2.0 \\
\midrule
Left shoulder  & 39.8 & 38.9 & 39.6 & 41.0 \\
Right shoulder & 39.3 & 38.1 & 38.6 & 40.2 \\
Left elbow     & 18.1 & 17.7 & 18.3 & 20.2 \\
Right elbow    & 17.5 & 16.7 & 18.2 & 20.2 \\
Left hip       & 29.6 & 29.0 & 29.9 & 32.5 \\
Right hip      & 32.0 & 31.7 & 32.8 & 36.7 \\
Left knee      & 16.1 & 16.6 & 17.7 & 17.8 \\
Right knee     & 18.9 & 18.7 & 19.6 & 20.6 \\
\bottomrule
\end{tabular}
\end{table}

\subsection{Inter-joint Error Correlation}

Tables~\ref{tab:correlation_theta} and \ref{tab:correlation_psi} list
high-correlation pairs ($|r|>0.7$) for $\Delta\theta$ and $\Delta\psi$,
respectively.

\begin{table}[!t]
\centering
\caption{High-correlation joint pairs in XY plane ($\Delta\theta$; $|r|>0.7$).
  All pairs: $p<0.001$ vs.\ $H_0:\rho=0$ (two-sided; $n>21{,}000$
  frame--camera observations). Approximate 95\% CIs for $\rho$ (Fisher $z$)
  are extremely narrow.}
\label{tab:correlation_theta}
\small
\begin{tabular}{llr}
\toprule
Joint 1        & Joint 2        & $r$ \\
\midrule
LEFT\_ELBOW    & RIGHT\_ELBOW   & $-0.862$ \\
LEFT\_ELBOW    & LEFT\_SHOULDER &  $0.788$ \\
RIGHT\_ANKLE   & RIGHT\_KNEE    &  $0.767$ \\
LEFT\_ANKLE    & LEFT\_KNEE     &  $0.766$ \\
RIGHT\_ELBOW   & RIGHT\_SHOULDER&  $0.710$ \\
\bottomrule
\end{tabular}
\end{table}

\begin{table}[!t]
\centering
\caption{High-correlation joint pairs in XZ plane ($\Delta\psi$; $|r|>0.7$).
  All pairs: $p<0.001$ vs.\ $H_0:\rho=0$ (two-sided; $n>21{,}000$
  frame--camera observations). Approximate 95\% CIs for $\rho$ (Fisher $z$)
  are extremely narrow.}
\label{tab:correlation_psi}
\small
\begin{tabular}{llr}
\toprule
Joint 1        & Joint 2         & $r$ \\
\midrule
LEFT\_HIP      & RIGHT\_HIP      & $-0.840$ \\
RIGHT\_ELBOW   & RIGHT\_SHOULDER &  $0.770$ \\
RIGHT\_ELBOW   & RIGHT\_WRIST    &  $0.769$ \\
LEFT\_ELBOW    & LEFT\_SHOULDER  &  $0.768$ \\
RIGHT\_SHOULDER& RIGHT\_WRIST    &  $0.726$ \\
LEFT\_ELBOW    & LEFT\_WRIST     &  $0.722$ \\
RIGHT\_ELBOW   & RIGHT\_HIP      &  $0.721$ \\
LEFT\_HIP      & RIGHT\_ELBOW    & $-0.705$ \\
\bottomrule
\end{tabular}
\end{table}

In the XY plane, the dominant feature is a strong negative correlation
between left and right elbow ($r=-0.862$), alongside positive same-side
elbow--shoulder and ankle--knee coupling. In the XZ plane, the hip
anti-correlation ($r=-0.84$) and upper-limb coupling ($r=0.72$--$0.77$) are
notable. Figure~\ref{fig:correlation_heatmaps} visualises the full
correlation matrices.

\begin{figure*}[!tb]
  \centering
  \begin{subfigure}[b]{0.49\textwidth}
    \centering
    \includegraphics[width=\linewidth]{figs/fig06a_corr_theta.png}
    \caption{$\Delta\theta$ (XY plane)}
  \end{subfigure}\hfill
  \begin{subfigure}[b]{0.49\textwidth}
    \centering
    \includegraphics[width=\linewidth]{figs/fig06b_corr_psi.png}
    \caption{$\Delta\psi$ (XZ plane)}
  \end{subfigure}
  \caption{Correlation matrices of inter-joint direction angle errors.
    (a)~$\Delta\theta$ corresponds to Table~\ref{tab:correlation_theta};
    (b)~$\Delta\psi$ corresponds to Table~\ref{tab:correlation_psi}.
    Source: \texttt{05\_direction\_detection/output/correlation\_analysis/}.}
  \label{fig:correlation_heatmaps}
\end{figure*}

\subsection{3-D Position Error (MPJPE)}

As a supplementary positional metric, we compute the Mean Per-Joint
Position Error (MPJPE), i.e.\ the Euclidean distance between estimated and
ground-truth 3-D landmark positions in MediaPipe world coordinate space
(units approximately equal to metres for a 1.7\,m subject).
The overall mean across all 12 joints is $0.363$\,m.
Hip joints serve as the coordinate origin and therefore show near-zero
error ($0.044$\,m), while distal limb joints exhibit larger values:
ankle ($0.71$\,m), wrist ($0.36$--$0.40$\,m), and elbow
($0.31$--$0.32$\,m).
The large ankle error primarily reflects depth-axis uncertainty in
monocular lifting: ankle depth is estimated relative to the hip origin
with no additional geometric constraint.
Positional errors also accumulate along the kinematic chain from the
hip-centred root toward distal joints, compounding distal uncertainty.
These values are consistent with the known difficulty of monocular 3-D
pose estimation and the absence of multi-view fusion in MediaPipe.

\subsection{Camera Viewpoint Dependency}
\label{subsec:camera_dep}

To characterise viewpoint dependence, we identified for each of the
107 frames the camera with the lowest and highest mean $|\Delta\theta|$
across joints (computed from
\path{05_direction_detection/output/processed_data/frame_camera_summary.csv},
107 frames $\times$ 201 cameras per frame).
The \emph{best} camera per frame averages $|\Delta\theta|=8.8^{\circ}$ over frames
(per-frame minima and maxima: $6.1^{\circ}$ and $15.8^{\circ}$).
The \emph{worst} camera averages $80.1^{\circ}$ (min $55.6^{\circ}$, max $148.1^{\circ}$),
i.e.\ roughly a nine-fold gap in mean $|\Delta\theta|$ relative to the best camera,
underscoring how strongly camera placement affects estimates.

% ============================================================
% 5. DISCUSSION  (target: ~2 pages)
% ============================================================
\section{Discussion}

\subsection{Coordinate System Unification Effect}

The Y-axis orientation mismatch between Unity and MediaPipe causes a
systematic sign error in $\theta$ for every joint. Because
$\theta_{\text{MP}} \approx -\theta_{\text{GT}}$ when the Y-axis is
inverted, the uncorrected direction angle error satisfies
$\Delta\theta \approx -2\theta_{\text{GT}}$; for a joint at
$\theta_{\text{GT}} \approx 60^{\circ}$ this yields $|\Delta\theta| \approx 120^{\circ}$.
After applying $y \leftarrow -y$, the reduction in mean $|\Delta\theta|$
for shoulders and elbows (and all other joints) is summarised in
Table~\ref{tab:yflip_effect}~\cite{mediapipe_coord_issue}.
Shoulder and ankle joints benefit the most ($\Delta|\Delta\theta|>155^{\circ}$),
because their ground-truth $\theta$ values are close to $\pm90^{\circ}$,
so the sign inversion nearly doubles the apparent error.
Wrist joints show the smallest improvement ($<10^{\circ}$) owing to
large estimation variance that partially masks the systematic sign error.

\begin{table}[!t]
\centering
\caption{Effect of Y-axis coordinate unification: mean $|\Delta\theta|$ (deg)
  before and after applying $y \leftarrow -y$ to MediaPipe coordinates.
  Source: \texttt{05\_direction\_detection} exports (before vs.\ after
  applying \texttt{coordinate\_transform.py}; see Sec.~3.4).}
\label{tab:yflip_effect}
\small
\begin{tabular}{lrrr}
\toprule
Joint          & Before & After & Reduction \\
\midrule
Left shoulder  & 169.5 & 10.4 & 159.1 \\
Right shoulder & 171.2 &  9.3 & 161.9 \\
Left elbow     & 121.6 & 58.3 &  63.3 \\
Right elbow    & 121.2 & 59.3 &  61.9 \\
Left wrist     &  94.3 & 84.3 &  10.0 \\
Right wrist    &  93.1 & 91.5 &   1.6 \\
Left knee      & 161.7 & 17.4 & 144.3 \\
Right knee     & 162.2 & 18.5 & 143.7 \\
Left ankle     & 168.7 & 11.3 & 157.4 \\
Right ankle    & 167.9 & 11.9 & 156.0 \\
\bottomrule
\end{tabular}
\end{table}

The residual $\approx58^{\circ}$ elbow error after unification reflects genuine
estimation bias (training-data viewpoint imbalance, arm-at-side prior),
not coordinate system mismatch.

\subsection{Systematic Bias in Depth Estimation}

MediaPipe's BlazePose estimates each landmark via heatmaps and regression,
outputting 3-D world coordinates relative to the hip
centre~\cite{blazepose}. The architecture does not explicitly incorporate
anatomical constraints such as bone-length consistency or pelvic
rigidity. The mean $|\Delta\theta|$ for elbow joints after coordinate
unification remains $\approx58^{\circ}$, indicating a persistent directional bias.

Many pose estimation models, including MediaPipe, are trained on corpora
such as COCO~\cite{coco} that are often described as skewing toward everyday,
near-frontal, standing- or walking-like poses, which can bias elbow direction
toward a downward-hanging posture. The opposite-sign pattern in left vs.\
right elbow errors ($r=-0.862$) is \emph{consistent with} left--right flip
data augmentation commonly used in pose-model training (we do not have access
to proprietary MediaPipe training details). The sample SD of elbow
$\Delta\theta$ ($\approx23$--$25^{\circ}$) is small relative to the mean
bias, suggesting that the residual systematic error is comparatively
predictable and amenable to post-processing.

\subsection{Left-Right Symmetric Hip Depth Error}

For hip joints, the mean $|\Delta\psi|\approx90^{\circ}$ and the strong negative
correlation ($r=-0.8402$) indicate that the depth-direction errors are
anti-correlated: when one hip is estimated as forward, the other is
estimated as backward, creating a phantom pelvic-rotation artefact. The
identical absolute-value heatmap patterns for left and right hip
(Figure~\ref{fig:hip_psi_heatmap}) are consistent with
$\Delta\psi_{\text{L}} \approx -\Delta\psi_{\text{R}}$ at each observation.

In MediaPipe's skeleton model the pelvis is represented as the segment
connecting landmark~23 (left hip) and landmark~24 (right hip). This linear
representation without a multi-segment rigid structure (e.g., sacrum and
ilium) imposes no rigidity constraint on the relative depth of the two
hips during 3-D lifting. Consequently, even a small horizontal pixel
difference in the 2-D image can be interpreted as placing one hip closer
and the other farther from the camera, producing the symmetric error
pattern. Model improvements that impose pelvis rigidity as a hard or soft
constraint are needed.

\subsection{Upper-Limb Coupled Error}

The high positive correlations among shoulder, elbow, and wrist in
$\Delta\psi$ ($r=0.72$--$0.77$) suggest that errors propagate
hierarchically along the kinematic chain, or that the upper limb is
implicitly treated as a rigid unit during 3-D lifting. Incorporating
anatomical bone-length constraints via graph neural networks or explicit
kinematic models could reduce this propagation.

\subsection{Viewpoint Dependence and Training-Data Bias}

The statistics in Section~\ref{subsec:camera_dep} (nine-fold mean
$|\Delta\theta|$ gap between best and worst cameras) and the elevated MAE at
Y\,=\,2.0\,m in Table~\ref{tab:mae_by_layer} jointly support the interpretation
that BlazePose training data (primarily COCO~\cite{coco} and fitness-style
recordings~\cite{blazepose}) tends to emphasise frontal and near-frontal
views, with lateral and top-down cameras relatively underrepresented, an
interpretation consistent with our layered MAE pattern.

\subsection{Practical Implications}

MediaPipe provides useful 2-D joint localisation but its 3-D depth
estimation is fundamentally limited by monocular ambiguity.
For applications where relative joint angles in the image plane suffice
(e.g., stroke-count analysis in swimming, gross-motor rehabilitation
screening) MediaPipe is adequate. For applications requiring 3-D body
orientation (e.g., gait analysis, ergonomic risk assessment of torso
rotation) the systematic biases identified here (in particular hip
$\Delta\psi$ and elbow $\Delta\theta$) must be compensated.

\paragraph{Deployment checklist (informal).}
\textbf{Camera placement:} prefer near-frontal, eye-level views when
3-D orientation matters; avoid relying on a single overhead (bird's-eye)
camera for quantitative joint-angle work.
\textbf{Views to treat cautiously:} large polar angles and Y$\approx$2.0\,m
layers showed elevated MAE (Table~\ref{tab:mae_by_layer}).
\textbf{Mandatory corrections:} apply $y \leftarrow -y$ (or equivalent)
before interpreting $\theta$ when mixing Unity (Y-up) GT with MediaPipe
world landmarks; plan compensations for hip $\Delta\psi$ and residual
elbow $\Delta\theta$ whenever pelvis orientation or forearm direction
drives the application logic.

\subsection{Limitations and external validity}

Limitations are grouped by how strongly they constrain interpretation:
\begin{itemize}\itemsep2pt
  \item \textbf{Avatar and appearance (high impact on transfer to humans):}
        Y-Bot with uniform materials (no soft-tissue or clothing
        artefacts). Conclusions on \emph{magnitude} of pixel-level jitter
        and segmentation errors may shift on real humans, whereas
        coordinate-convention and monocular-depth issues are expected to
        persist qualitatively.
  \item \textbf{Single motion clip (moderate impact):}
        One looped walk limits generalisation to sports, sitting, or
        fast articulation; viewpoint findings should be re-checked per
        motion domain.
  \item \textbf{Lighting and domain gap (moderate for appearance-driven
        failure modes):}
        Studio lighting reduces shadow/occlusion extremes; failure rates
        under hard backlighting or motion blur are not measured here.
  \item \textbf{Single estimator (scope boundary):}
        Numbers are MediaPipe-specific; alternative monocular estimators
        need their own calibration but can reuse the evaluation protocol
        (Sec.~3.4 and Table~\ref{tab:repro_pipeline}).
\end{itemize}

% ============================================================
% 6. CONCLUSION  (target: ~0.5 page)
% ============================================================
\section{Conclusion}

Unity simulation with 505 multi-view cameras quantified systematic biases in
MediaPipe pose estimation. Joint MAE and layer-wise viewpoint dependence are
summarised in Tables~\ref{tab:joint_angle_error} and~\ref{tab:mae_by_layer};
the effect of $y \leftarrow -y$ appears in Table~\ref{tab:yflip_effect};
interpretation of hip $\Delta\psi$ anti-correlation, upper-limb coupling, and
residual elbow error is given in Section~5 and
Tables~\ref{tab:correlation_theta}--\ref{tab:correlation_psi}.
The reported statistics describe behaviour under this synthetic, reproducible
GT setup; extrapolation to real subjects and in-the-wild video requires
separate validation. The quantification nevertheless supplies actionable
structure (viewpoint imbalance, coordinate conventions) for deployment design
and for ancillary, bias-aware use of pose outputs.

Future work should address pelvic rigidity constraints, learning from
multi-view data, and porting the protocol to other estimators.

% ============================================================

\acknowledgment

This work was supported by the Japan Science and Technology Agency (JST)
under the Next-Generation Human Resource Development Program (Global Science
Campus Initiative).
This work was supported in part by the Japan Society for the Promotion of
Science (JSPS) Grants-in-Aid for Scientific Research (Grant Numbers:
23K25723, Principal Investigator: Koji Sugio; and 25K06208, Principal
Investigator: Tsukasa Kato).


\begin{thebibliography}{99}

\bibitem{zenodo_mp_csv}
F.~Taira: ``MediaPipe processed pose CSVs'', Zenodo, dataset,
\url{https://doi.org/10.5281/zenodo.19296530} (Mar.~2026)

% --- Original 3 references (kept from paper_en.md) ----------
\bibitem{mediapipe}
Google LLC: ``MediaPipe Pose Landmarker'',
\url{https://ai.google.dev/edge/mediapipe/solutions/vision/pose_landmarker}
(2024)

\bibitem{mediapipe_coord_issue}
Google: ``Orientation of pose world landmarks'', GitHub Issue \#3370,
MediaPipe repository,
\url{https://github.com/google/mediapipe/issues/3370} (2022)

\bibitem{blazepose}
V.~Bazarevsky, I.~Grishchenko, K.~Raveendran, T.~Zhu, F.~Zhang,
and M.~Grundmann: ``BlazePose: On-device real-time body pose tracking'',
\textit{arXiv preprint arXiv:2006.10204} (2020)

\bibitem{unity}
Unity Technologies: ``HumanBodyBones'', \textit{Unity Scripting Reference},
\url{https://docs.unity3d.com/ScriptReference/HumanBodyBones.html} (2024)

\bibitem{coco}
T.-Y.~Lin, M.~Maire, S.~Belongie, J.~Hays, P.~Perona, D.~Ramanan,
P.~Dollár, and C.~L.~Zitnick:
``Microsoft COCO: Common Objects in Context'',
in \textit{Proc. ECCV}, pp.~740--755 (2014)

\bibitem{openpose}
Z.~Cao, G.~Hidalgo, T.~Simon, S.-E.~Wei, and Y.~Sheikh:
``OpenPose: Realtime Multi-Person 2D Pose Estimation Using Part Affinity
Fields'', \textit{IEEE Trans. Pattern Anal. Mach. Intell.}, Vol.~43,
No.~1, pp.~172--186 (2021)

% --- Synthetic / simulation-based GT datasets ----------------------
\bibitem{surreal}
G.~Varol, J.~Romero, X.~Martin, N.~Mahmood, M.~J.~Black, I.~Laptev,
and C.~Schmid: ``Learning from synthetic humans'',
in \textit{Proc. CVPR}, pp.~109--117 (2017)

\bibitem{agora}
P.~Patel, C.-H.~P.~Huang, J.~Tesch, D.~T.~Hoffmann, S.~Tripathi,
and M.~J.~Black: ``AGORA: Avatars in geography optimized for regression
analysis'', in \textit{Proc. CVPR} (2021)

\bibitem{peoplesanspeople}
S.~E.~Ebadi, Y.-C.~Jhang, A.~Zook, S.~Dhakad, A.~Crespi, P.~Parisi,
S.~Borkman, J.~Hogins, and S.~Ganguly:
``PeopleSansPeople: A synthetic data generator for human-centric computer
vision'', \textit{arXiv preprint arXiv:2112.09290} (2021)

% --- Monocular 3D pose estimation / depth ambiguity ----------------
\bibitem{mehta2017}
D.~Mehta, H.~Rhodin, D.~Casas, P.~Fua, O.~Sotnychenko, W.~Xu,
and C.~Theobalt: ``Monocular 3D human pose estimation in the wild using
improved CNN supervision'', in \textit{Proc. Int.\ Conf.\ 3D Vision (3DV)},
pp.~506--516, \url{https://doi.org/10.1109/3DV.2017.00064} (2017)

\bibitem{vnect}
D.~Mehta, S.~Sridhar, O.~Sotnychenko, H.~Rhodin, M.~Shafiei,
H.-P.~Seidel, W.~Xu, D.~Casas, and C.~Theobalt:
``VNect: Real-time 3D human pose estimation with a single RGB camera'',
\textit{ACM Trans.\ Graph.}, Vol.~36, No.~4, pp.~44:1--44:14 (2017)

% --- Skeletal constraint / transformer-based models ----------------
\bibitem{poseformer}
C.~Zheng, S.~Zhu, M.~Mendieta, T.~Yang, C.~Chen, and Z.~Ding:
``3D human pose estimation with spatial and temporal transformers'',
in \textit{Proc. ICCV}, pp.~11656--11666 (2021)

\bibitem{motionbert}
W.~Zhu, X.~Ma, Z.~Liu, L.~Liu, W.~Wu, and Y.~Wang:
``MotionBERT: A unified perspective on learning human motion
representations'', in \textit{Proc. ICCV} (2023)

\end{thebibliography}


\begin{biography}
\profile{n}{Fumimaro Taira}{%
Second-year student, Okinawa Prefectural Kaiko High School.
Two-stage student at Ryudai Kagaku-in under the University of the Ryukyus
SEARCH programme, affiliated with Prof.~Tsukasa Kato's laboratory;
engaged in motion-tracking research.}
\profile{m}{Tsukasa Kato}{%
Received the B.S. degree from Tokyo University of Fisheries
(now Tokyo University of Marine Science and Technology) in March 1999.
While serving at Okinawa Prefectural high schools, he completed the
graduate program in the Department of Intelligent Information Engineering,
University of the Ryukyus, in March 2021, and received the Ph.D. degree in
engineering. He is currently an Associate Professor in the Professional
Degree Program for Advanced Teacher Education, Graduate School of Education,
University of the Ryukyus. His research interests include the development of
video-based lessons and mobile learning systems, the development of
technical instructional videos and evaluation of their effectiveness, and
the educational use and application of generative AI.}
\profile{n}{Jin Afuso}{%
Received the B.E. degree in information engineering from the University of
the Ryukyus in March 2004. He completed the doctoral program in the Major
of Intelligent Information Engineering, Graduate School of Science and
Engineering, University of the Ryukyus, in March 2013, and received the
Ph.D. degree in engineering. In April 2013, he joined the Knowledge Media
Laboratory, Hokkaido University, as a Project Researcher. In April 2015, he
became an Assistant Professor by Special Appointment at the University of
the Ryukyus Hospital. After serving from December 2017 as a Project
Researcher in the Social Infrastructure Program, Faculty of Engineering,
University of the Ryukyus, he is currently engaged in artificial
intelligence research as a data scientist.}
\end{biography}

\end{document}
